% 第11章第3节：从树到一般图（环带来的接口变化、四环算例、一般图动态规划的构建目标）。

\section{从树到一般图}
\label{sec:graph-dp-tree-to-general}

上一节的每一步都依赖一个事实：删去依赖树的一条边，变量就分成两侧，两侧只通过一个变量相连。依赖图含环时，这个事实不再成立。本节先说明环怎样改变接口，再用一个四环算例把变化落到具体数值上，最后说明一般图上的动态规划要事先构建什么。

\subsection{环带来的接口变化}
\label{sec:graph-dp-cycle-interface}

考虑四个变量排成的环$1-2-3-4-1$，每条边上有一个成对成本。删去边$\{1,2\}$后，剩下的是路径$2-3-4-1$，仍然连通：$x_1$不直接与$x_2$相连，却可以经由$x_4$和$x_3$影响$x_2$的最优选择。因此，沿一条边传递一张只以$x_2$为索引的报价单，无法概括$x_1$一侧，因为根本不存在只经过$x_2$与其余部分相连的“$x_1$一侧”。

链和树上的做法可以换一种说法，使它不依赖“一侧”。叶顶点$u$只出现在$\phi_u$和$\phi_{u,p(u)}$两项中；计算消息$m_{u\to p(u)}$，就是把这两项相加后对$x_u$取最小，然后用得到的表替换这两项。这个动作对任意局部函数集合都有意义。

设当前有一组函数，要汇总变量$x_i$。先找出当前全部作用域含$i$的函数，记为$\mathcal F_i$；它们的作用域之并去掉$i$，记为$S_i$。不含$i$的函数与$x_i$无关，可以留在最小化之外；含$i$的函数则必须一起处理。于是对$\bm x_{S_i}$的每种取值，把$\mathcal F_i$中的函数相加，再对$x_i$取最小，得到一个新函数
\begin{equation}
\psi_i(\bm x_{S_i})=\min_{x_i\in\mathcal X_i}\sum_{f\in\mathcal F_i}f .
\label{eq:graph-dp-elimination-step}
\end{equation}
然后从当前函数集合中删去$\mathcal F_i$，加入$\psi_i$；不含$i$的函数保持原样。这一操作称为消去$x_i$，按某个次序逐个消去变量的方法称为\term{变量消元}（variable elimination）。

消去$x_i$可以理解为替$x_i$提前做好决定。$x_i$的好坏只取决于$S_i$中那些与它同在某个函数里的变量；对这些变量的每一种取值，先算出$x_i$最好取什么、最少付多少，写成一张以$\bm x_{S_i}$为索引的表，此后其余变量只需查这张表，不必再关心$x_i$。由式~\eqref{eq:graph-dp-min-rules}的前两条，不含$i$的函数可以移出对$x_i$的最小化，所以替换前后，剩余变量每种取值下的最小成本不变；第~\ref{sec:graph-dp-min-cost-correctness}小节给出完整证明。

在树上消去叶顶点$u$时，$S_u=\{p(u)\}$，新函数就是消息$m_{u\to p(u)}$，接口仍是单个变量。在四环上消去$x_1$，情况不同：含$1$的函数是$\phi_{1,2}$与$\phi_{1,4}$，所以$S_1=\{2,4\}$，新函数$\psi_1(x_2,x_4)$同时读取$x_2$和$x_4$。在原来的环上，$2$与$4$并不相邻；消去$x_1$以后，它们通过$\psi_1$直接耦合在一起。

能否不要这张二元表，而给$2$和$4$各发一张单变量表？分别计算
\[
g(x_2)=\min_{x_1}\phi_{1,2}(x_1,x_2),\qquad h(x_4)=\min_{x_1}\phi_{1,4}(x_1,x_4)
\]
相当于允许两项各自为$x_1$挑选取值。由于每个$x_1$都满足$\phi_{1,2}+\phi_{1,4}\geq g(x_2)+h(x_4)$，对$x_1$取最小后总有
\[
g(x_2)+h(x_4)\leq\psi_1(x_2,x_4),
\]
但两边一般不相等：$x_1$只有一个，它在两项中必须取同一个值。下一小节的数值会说明，这个差异不仅是这一种拆法的缺陷，而且任何单变量表之和都无法代替$\psi_1$。

还有一种看似自然的做法，是在含环的依赖图上照搬式~\eqref{eq:graph-dp-tree-message}，让各顶点反复交换消息。这时同一部分成本会沿环被重复计入，上一节的正确性证明不再适用。这样的迭代是否收敛、收敛到什么，都需要另外分析；本部分只讨论精确计算。

\subsection{四环算例}
\label{sec:graph-dp-four-cycle}

本小节固定一个具体问题，后两章沿用它，分别演示两种结构构建和统一求值，避免更换例子掩盖它们之间的对应关系。四个二值变量$x_1,\ldots,x_4$排成环$1-2-3-4-1$，四条边上各有一个成对成本，没有一元项：
\begin{equation}
F(\bm x)=\phi_{1,2}(x_1,x_2)+\phi_{2,3}(x_2,x_3)+\phi_{3,4}(x_3,x_4)+\phi_{1,4}(x_1,x_4).
\label{eq:graph-dp-four-cycle}
\end{equation}
四张表由表~\ref{tab:graph-dp-four-cycle-costs}给出，每张表的行对应作用域中编号较小的变量。其中两个$+\infty$表示禁止的组合：$\phi_{3,4}(1,1)=+\infty$禁止$x_3,x_4$同时取$1$，$\phi_{1,4}(0,1)=+\infty$禁止$x_1=0$而$x_4=1$。

\begin{table}[!htb]
  \centering
  \begin{tabular}{cccc}
    \toprule
    $\phi_{1,2}(x_1,x_2)$ & $\phi_{2,3}(x_2,x_3)$ & $\phi_{3,4}(x_3,x_4)$ & $\phi_{1,4}(x_1,x_4)$\\
    \midrule
    $\begin{bmatrix}3&0\\1&2\end{bmatrix}$ &
    $\begin{bmatrix}1&3\\1&1\end{bmatrix}$ &
    $\begin{bmatrix}0&1\\0&+\infty\end{bmatrix}$ &
    $\begin{bmatrix}3&+\infty\\3&0\end{bmatrix}$\\
    \bottomrule
  \end{tabular}
  \caption{四环算例的局部成本。每张表的行对应作用域中编号较小的变量取$0,1$，列对应编号较大的变量取$0,1$；例如$\phi_{1,4}(0,1)=+\infty$。}
  \label{tab:graph-dp-four-cycle-costs}
\end{table}

\paragraph{枚举基准}
先用枚举给出核验基准。表~\ref{tab:graph-dp-four-cycle-enumeration}按字典序列出全部$16$个配置的总成本。其中$6$个配置被禁止：$x_4=1$而$x_1=0$的四个，以及$x_3=x_4=1$、$x_1=1$的两个。其余$10$个可行。

最小成本为$F^\star=3$，唯一的最优配置是$\bm x^\star=(1,0,0,1)$，此时四项依次为
\[
\phi_{1,2}(1,0)=1,\quad \phi_{2,3}(0,0)=1,\quad \phi_{3,4}(0,1)=1,\quad \phi_{1,4}(1,1)=0.
\]

\begin{table}[!htb]
  \centering
  \small
  \begin{tabular}{@{}cccccccc@{}}
    \toprule
    $0000$ & $0001$ & $0010$ & $0011$ & $0100$ & $0101$ & $0110$ & $0111$\\
    $7$ & $+\infty$ & $9$ & $+\infty$ & $4$ & $+\infty$ & $4$ & $+\infty$\\
    \midrule
    $1000$ & $1001$ & $1010$ & $1011$ & $1100$ & $1101$ & $1110$ & $1111$\\
    $5$ & $3$ & $7$ & $+\infty$ & $6$ & $4$ & $6$ & $+\infty$\\
    \bottomrule
  \end{tabular}
  \caption{四环算例全部配置的总成本。配置按$(x_1,x_2,x_3,x_4)$的字典序排列，前两行是$x_1=0$的八个配置，后两行是$x_1=1$的八个配置。}
  \label{tab:graph-dp-four-cycle-enumeration}
\end{table}

\paragraph{消去$x_1$}
按式~\eqref{eq:graph-dp-elimination-step}消去$x_1$。对$(x_2,x_4)$的每种取值，比较$x_1=0,1$两个候选$\phi_{1,2}(x_1,x_2)+\phi_{1,4}(x_1,x_4)$：
\[
\begin{aligned}
\psi_1(0,0)&=\min(3+3,\ 1+3)=4,&\qquad \psi_1(0,1)&=\min(3+\infty,\ 1+0)=1,\\
\psi_1(1,0)&=\min(0+3,\ 2+3)=3,&\qquad \psi_1(1,1)&=\min(0+\infty,\ 2+0)=2,
\end{aligned}
\]
达到最小值的$x_1$依次为$1,1,0,1$。

与之比较，分别取最小的单变量表是
\[
\begin{aligned}
g(x_2)&=\min_{x_1}\phi_{1,2}(x_1,x_2)=(1,0),\\
h(x_4)&=\min_{x_1}\phi_{1,4}(x_1,x_4)=(3,0).
\end{aligned}
\]
它们的和在$(x_2,x_4)=(1,1)$处为$0$，而$\psi_1(1,1)=2$：$g(1)=0$要求$x_1=0$，$h(1)=0$要求$x_1=1$，两项选的$x_1$互相矛盾。

更进一步，任何形如$g'(x_2)+h'(x_4)$的有限函数都满足
\[
\bigl[g'(0)+h'(0)\bigr]+\bigl[g'(1)+h'(1)\bigr]=\bigl[g'(0)+h'(1)\bigr]+\bigl[g'(1)+h'(0)\bigr],
\]
而$\psi_1$不满足这一等式：
\[
\psi_1(0,0)+\psi_1(1,1)=6\neq4=\psi_1(0,1)+\psi_1(1,0).
\]
所以$\psi_1$不能拆成两张单变量表之和，消去$x_1$以后的接口必须是$\{2,4\}$这两个变量的联合取值。

\paragraph{完成求值}
消去$x_1$后，剩下的函数是$\psi_1(x_2,x_4)$、$\phi_{2,3}(x_2,x_3)$与$\phi_{3,4}(x_3,x_4)$，只涉及$x_2,x_3,x_4$三个变量。把三项相加，按$(x_2,x_3,x_4)$的字典序得到
\[
(5,\ 3,\ 7,\ +\infty,\ 4,\ 4,\ 4,\ +\infty),
\]
最小值$3$在$(x_2,x_3,x_4)=(0,0,1)$处唯一达到。

再查$x_1$的选择：$(x_2,x_4)=(0,1)$时$x_1=1$。恢复的配置$(1,0,0,1)$与枚举结果一致。

这里同时处理的变量最多是三个：消去$x_1$时要看$\{1,2,4\}$的全部$8$种取值，最后又要看$\{2,3,4\}$的$8$种取值。第~\ref{chap:general-graph-dp-computation}章会把这一计算写成统一的递推，并用同一组数据检验最小成本、保留变量的成本表、计数和加权求和。

\subsection{一般图动态规划的构建目标}
\label{sec:graph-dp-construction-goal}

四环算例中的计算可以看成两个局部子问题。第一个子问题负责变量组$\{1,2,4\}$，固定接口$(x_2,x_4)$后汇总内部变量$x_1$，输出表$\psi_1$；第二个子问题负责$\{2,3,4\}$，读入$\psi_1$和其余两项，汇总全部剩余变量。链和树上的递推也是这样一族子问题，只是每个接口恰好是单个变量。

一般图上的动态规划要先构建这样一族子问题，它必须满足三项要求：
\begin{enumerate}
  \item 每个局部项都放进某个子问题，该子问题处理的变量包含它的整个作用域，并且每项只计一次；
  \item 摘要足够：一个子问题的内部变量与尚未处理部分的全部耦合都经过接口，所以以接口为索引的表不会丢失信息；
  \item 依赖可以排序：每个子问题只使用此前已完成的子问题的输出，最后汇总成整体结果。
\end{enumerate}
这三项都只涉及作用域和变量组，不涉及函数的具体数值，所以构建可以在求值之前单独完成。

树上递推正是满足这些要求的特例。把依赖树的一元项和成对项作为局部函数，按后序消去非根顶点：轮到$u$时，它的子点都已消去，子点留下的新函数与$\phi_u$、$\phi_{u,p(u)}$是仅有的含$u$的函数，所以$S_u=\{p(u)\}$，新函数$\psi_u$就是消息$m_{u\to p(u)}$，达到最小值的$x_u$就是回溯表$b_{u\to p(u)}$。

构建这族子问题有两种描述方式。变量消元用一个处理次序生成它们：每一步处理的变量组是当前变量及其$S_i$，新函数交给后面第一个用到它的步骤。树分解则把变量组（称为袋）和它们之间的接口显式组织成一棵树。第~\ref{chap:general-graph-dp-construction}章说明这两种构建怎样互相转换、怎样衡量其规模、怎样构造；第~\ref{chap:general-graph-dp-computation}章在构建好的结构上完成数值计算。四环算例已经显示，构建的质量体现为最多同时处理多少个变量，这里是三个。
